% ================================================================
\section{Related Work}
% ================================================================

\subsection{Open Code Language Models}

Code generation has become a standard benchmark for large language
models. Building on the decoder-only Transformer~\cite{Vaswani2017Attention},
a succession of open models has narrowed the gap to proprietary
systems. StarCoder~\cite{Li2023StarCoder} trains on the permissively
licensed Stack corpus with multi-query attention and a
fill-in-the-middle objective. Code Llama~\cite{Roziere2023CodeLlama}
extends Llama~2 with a long-context fine-tuning stage and infilling
support. DeepSeek-Coder~\cite{Guo2024DeepSeek} trains on a
project-level corpus with an infilling objective, achieving
state-of-the-art results at several scales. Qwen2.5-Coder
\cite{Hui2024Qwen25} reports leading scores among open models at the
1.5B parameter scale we adopt.

These efforts target general-purpose, large-scale code generation.
Completus differs in goal: rather than pretraining a new model, we
specialise an existing small model for a narrow, latency-sensitive use
case where size and local runnability matter more than peak benchmark
performance.

\subsection{Parameter-Efficient Fine-Tuning for Code}

Low-rank adaptation (LoRA)~\cite{Hu2021LoRA} freezes pretrained
weights and injects trainable rank-decomposition matrices, cutting
trainable parameters by orders of magnitude. QLoRA~\cite{Dettmers2023QLoRA}
extends this by quantising the frozen base model to 4-bit NF4, making
gradient computation feasible on a single consumer GPU. AdaLoRA
\cite{Zhang2023AdaLoRA} refines rank allocation by assigning capacity
adaptively according to weight-matrix importance scores.

Closest to our work, Weyssow et al.~\cite{Weyssow2023PEFT} evaluate
PEFT methods specifically for code generation and confirm that LoRA
adapters recover most of full fine-tuning quality at a fraction of
the cost. The comprehensive PEFT survey by Han et
al.~\cite{Han2024PEFTSurvey} situates LoRA within the broader family
of additive, selective, and reparameterisation methods.

\subsection{Training Data Curation for Code}

Xue et al.~\cite{Xue2025CleanCode} demonstrate that removing code
smells from the training set improves downstream model quality. Lee et
al.~\cite{Lee2021Dedup} show that near-deduplication reduces
memorisation and train-test leakage while cutting training cost.
Bavarian et al.~\cite{Bavarian2022FIM} show that models learn
infilling at no cost to left-to-right generation through a simple data
transformation using sentinel tokens. WizardCoder
\cite{Luo2023WizardCoder} demonstrates that instruction-style
fine-tuning further lifts code generation quality.

\subsection{Local Serving and Quantisation}

Kurt~\cite{Kurt2026Quant} provides a unified evaluation of llama.cpp
GGUF quantisation formats across size, speed, and quality
trade-offs, guiding our choice of the q4\_k\_m build. Whereas code
models are typically evaluated as cloud-hosted endpoints, Completus
targets a fully local serving workflow.
